\section{Wire Format Engine}

\subsection{Encoding}
The serialization engine in \texttt{ctpb::wire} manipulates the byte stream directly, without intermediate allocation. Encoding proceeds in two passes: \texttt{byte\_size()} computes the exact wire length from compile-time field metadata, the output buffer is allocated once, and the fields are then written through a raw pointer with no reallocation or virtual dispatch. Varints emit seven payload bits per byte with the high bit signalling continuation; signed integers are Zigzag-converted with single-instruction shifts and XORs; and repeated numeric scalars, which proto3 packs by default, have their payload size computed inline and are written as one contiguous length-delimited record under a single tag.

\subsection{Deserialization and Error Recovery}
The \texttt{ctpb::wire::reader} operates on a non-owning \texttt{std::span<const std::byte>} and bounds-checks every read: varint decoding refuses to advance past the buffer end or beyond the ten-byte maximum width, and length-delimited reads validate the declared length against the remaining bytes before yielding a subspan. Failures are reported through a typed \texttt{decode\_error} rather than an exception, so decoding remains usable in \texttt{noexcept} contexts.

Unknown field numbers are not an error. The \texttt{skip\_field} method uses the wire type alone to advance past an unrecognized integer, fixed-width block, or length-delimited payload, preserving proto3's forward-compatibility guarantee when a peer sends a newer schema version.
